\textbf{LLM as Agent}
% Traditional Reinforcement Learning provides general solutions for decision-making but lacks sample efficiency and generalization \citep{pourchot2018cem}. The emergent reasoning and instruction-following abilities of LLMs \citep{wei2022emergent} enable them to become proficient agents, excelling in zero-shot generalization~\citep{yao2022react, autogpt, wang2023voyager}. The central method for employing LLMs as agents is to prompt them with task instructions and environmental context, enabling them to produce actionable responses~\citep{autogpt, xie2023openagents}. Other methods involve specialized training to repurpose LLMs into adept agents \citep{xu2023lemur, reed2022generalist, driess2023palm}. We benchmark both general \citep{openai2023gpt, touvron2023llama, chiang2023vicuna} and agent LLMs \citep{xu2023lemur} to study LLMs as agents. Meanwhile, several work addresses dimensional aspects of agentic abilities, with a focus on the ability to ground goals to executable actions \citep{gu2022don, ahn2022can}, the ability to model the world \citep{lecun2022path}, the ability to perform step-by-step planning \citep{song2023llm}, and self-reflection ability \citep{madaan2023self, wang2023mint}. Examining various agentic skills is crucial to fully understand the advantages and limitations of LLMs as agents.
Traditional Reinforcement Learning offers general decision-making solutions but struggles with sample efficiency and generalization \citep{pourchot2018cem}. In contrast, the emergent reasoning and instruction-following abilities of LLMs \citep{wei2022emergent} enable them to excel as agents~\citep{yao2022react, autogpt, wang2023voyager}. The primary method for employing LLMs as agents involves prompting them with task instructions and environmental context to generate actionable responses~\citep{autogpt, xie2023openagents}. Specialized training can further enhance their agentic capabilities \citep{xu2023lemur, reed2022generalist, driess2023palm}. We benchmark both general \citep{openai2023gpt, touvron2023llama, chiang2023vicuna} and agent-specific LLMs \citep{xu2023lemur} to study their effectiveness as agents. Additionally, research explores various dimensions of agent abilities, including grounding goals to actions \citep{gu2022don, ahn2022can}, world modeling \citep{lecun2022path}, step-by-step planning \citep{song2023llm}, and self-reflection \citep{madaan2023self, wang2023mint}. Evaluating these skills is essential to understand limitations of LLMs as agents.

% On the other hand, agent ability is a critical feature of Large Language Models (LLMs), empowering them to actuate real-world changes \citep{ahn2022can, boiko2023autonomous,yang2023appagent} and acquire knowledge \citep{nakano2021webgpt}. Agent ability in LLMs is also thought to correlate with their chain-of-thought and reasoning skills \citep{zhang2023igniting}. Therefore, assessing LLMs as agents is essential for understanding LLMs generalist ability and driving their progress.

%\jh{cite AgentLM from Zhipu in this paragraph} (cited in sentence other methods involve specialized training...)

\textbf{Evaluating LLM in Decision Making Problems}
Several benchmarks and toolkits for LLM agents have been established, focusing on various tasks such as web-browsing, games, and tool use \citep{yao2022webshop, zhou2023webarena, shridhar2020alfworld, qin2023tool, wang2023voyager, ye2024tooleyes, kinniment2023evaluating}. A few other benchmarks provide a proof-of-concept study on specific LLM features, with \citet{wang2023mint} focusing on model interaction ability, and \citet{liu2023bolaa} examining agent structures. Recent works by \citet{liu2023agentbench, wu2023smartplay, mialon2023gaia} present a generalist challenge for LLM agents, please refer to Table \ref{table: dataset_comparison} for a comparison. 
% It is the first to offer a comprehensive evaluation framework to assess the capabilities and limitations of LLM agents and to offer insights into the fundamental attributes of LLMs. 
Note that recent progress in multimodal LLMs has spurred research into multimodal LLM agents \citep{Zheng2024GPT4VisionIA, yang2023appagent}. Our study focuses exclusively on text-based environments to assess LLM agent abilities via textual reasoning and actions in-depth.



% Despite our work share similar designs in terms of measuring generalist agent ability in unified and challenging 




